\documentclass[12pt,a4paper,oneside]{article}
\usepackage[brazil]{babel}
\usepackage[utf8]{inputenc}
\usepackage[dvipsnames]{xcolor}
\usepackage{listings}
\usepackage{olymp}

\setcounter{problem}{0}

\begin{document}

\begin{problem}{Ponta 
de flecha}{ponta}{2}


\begin{minipage}{.8\textwidth}
Na álgebra linear, uma matriz ponta de flecha é uma matriz quadrada que contém zeros em todas as entradas, exceto na primeira linha, primeira coluna e diagonal principal. Assim, ela tem o formato mostrado ao lado (considere que as entradas $*$ são números $\neq 0$).
\end{minipage}
\hfill
\begin{minipage}{.15\textwidth}
{
\small
\begin{bmatrix}
* & * & * & * & *\\
* & * & 0 & 0 & 0\\
* & 0 & * & 0 & 0\\
* & 0 & 0 & * & 0\\
* & 0 & 0 & 0 & *\\
\end{bmatrix}\\
}
\end{minipage}

%\Notes
%
%Observações, se tiver.

\InputFile

A primeira linha da entrada contém um número inteiro $N$, indicando a ordem da matriz. As $N$ linhas seguintes contêm exatamente $N$ números cada, indicando os valores das entradas da matriz. Restrições: $3 \leq N \leq 40$ e todos os demais valores inteiros entre $-1000$ e $1000$.

\OutputFile

Seu programa deve escrever ``\texttt{SIM}'' ou ``\texttt{NAO}'' na saída, informando se a matriz da entrada é ou não uma matriz ponta de flecha conforme definição acima.

% \Notes
% Preste atenção aos tipos das variáveis, pois $F(90) \approx 2^{61}$.

\Example

\begin{example}
\exmp{
5
1 2 3 4 5
2 4 0 0 0
3 0 6 0 0
4 0 0 8 0
5 0 0 0 1
}{
SIM
}%
\end{example}

\begin{example}
\exmp{
4
1 1 1 1
1 1 0 0
1 0 1 0
1 0 0 0
}{
NAO
}%
\end{example}

\begin{example}
\exmp{
4
1 1 1 1
1 1 0 0
1 0 1 0
1 1 0 1
}{
NAO
}%
\end{example}

\begin{example}
\exmp{
3
-1 5 7
8 4 0
2 0 -3
}{
SIM
}%
\end{example}

% \begin{example}
% \exmp{
% 4
% -9 11 -8 20
% 10 13 0 0 0
% -7 0 99 0 0
% -2 0 0 -4 0
% 18 0 0 0 20
% }{
% SIM
% }%
% \end{example}

%Comentários sobre o exemplo, se necessário.

\end{problem}

\end{document}
